\documentclass[10pt]{article}
% Shared LaTeX style for MIT 18.06 exam revision notes.
% Compile topic files with Tectonic, XeLaTeX, or LuaLaTeX.

\usepackage[a4paper,margin=0.72in]{geometry}
\usepackage{fontspec}
\usepackage{amsmath,amssymb,mathtools}
\usepackage{booktabs,array,tabularx}
\usepackage{enumitem}
\usepackage{titlesec}
\usepackage{fancyhdr}
\usepackage{microtype}
\usepackage{xcolor}

\setmainfont{Times New Roman}
\setsansfont{Arial}

\setlength{\parindent}{0pt}
\setlength{\parskip}{3.6pt}
\setlength{\abovedisplayskip}{6pt}
\setlength{\belowdisplayskip}{6pt}
\setlength{\abovedisplayshortskip}{4pt}
\setlength{\belowdisplayshortskip}{4pt}
\renewcommand{\arraystretch}{1.08}
\setlength{\tabcolsep}{4.5pt}

\titleformat{\section}
  {\normalfont\large\bfseries}
  {\thesection}
  {0.55em}
  {}
\titlespacing*{\section}{0pt}{10pt}{3pt}

\titleformat{\subsection}
  {\normalfont\normalsize\bfseries}
  {\thesubsection}
  {0.5em}
  {}
\titlespacing*{\subsection}{0pt}{7pt}{2pt}

\setlist[itemize]{leftmargin=1.35em,itemsep=1pt,topsep=2pt,parsep=0pt}
\setlist[enumerate]{leftmargin=1.55em,itemsep=1pt,topsep=2pt,parsep=0pt}

\pagestyle{fancy}
\fancyhf{}
\fancyfoot[C]{\scriptsize\color{black!45}MIT 18.06 Linear Algebra Exam Notes --- \thepage}
\renewcommand{\headrulewidth}{0pt}
\renewcommand{\footrulewidth}{0pt}

\newcommand{\notesubtitle}[1]{%
  {\centering\itshape #1\par}
  \vspace{2pt}
}

\newcommand{\source}[1]{%
  {\centering\small #1\par}
  \vspace{6pt}
}

\newcommand{\labelnote}[2]{%
  \par\smallskip
  \noindent\textbf{#1.}\quad #2
  \par\smallskip
}

\newcommand{\tightlabel}[1]{%
  \par\smallskip
  \noindent\textbf{#1.}\quad
}

\newcolumntype{Y}{>{\raggedright\arraybackslash}X}
\newcolumntype{L}[1]{>{\raggedright\arraybackslash}p{#1}}

\newenvironment{notetable}
  {\small\begin{tabularx}{\linewidth}}
  {\end{tabularx}}


\begin{document}

\begin{center}
{\LARGE 7.2 Singular Value Decomposition\par}
\end{center}
\notesubtitle{Exam notes for \(A=U\Sigma V^{T}\), singular values, and the four subspaces.}
\source{Based on handwritten MIT 18.06 revision notes.}

\section{Core Idea}

\labelnote{Core idea}{The SVD writes any matrix as orthogonal input directions, nonnegative stretching, and orthogonal output directions.}

For an \(m\times n\) matrix,
\[
A=U\Sigma V^{T}.
\]
Here \(V\) gives orthonormal directions in \(\mathbb{R}^n\), \(U\) gives orthonormal directions in \(\mathbb{R}^m\), and \(\Sigma\) contains the singular values.

\section{What the SVD Says}

\[
Av_i=\sigma_i u_i.
\]

This means: the right singular vector \(v_i\) is sent to the left singular vector \(u_i\), scaled by \(\sigma_i\ge 0\).

\begin{center}
\small
\begin{tabularx}{\linewidth}{@{}L{0.22\linewidth}Y Y@{}}
\toprule
Object & Lives in & Meaning \\
\midrule
\(v_i\) & \(\mathbb{R}^n\) & Orthonormal input direction. \\
\(u_i\) & \(\mathbb{R}^m\) & Orthonormal output direction. \\
\(\sigma_i\) & Nonnegative number & Stretch factor from \(v_i\) to \(u_i\). \\
\(\Sigma\) & \(m\times n\) diagonal shape & Stores singular values in decreasing order. \\
\bottomrule
\end{tabularx}
\end{center}

\section{How to Compute It}

The right singular vectors come from the symmetric matrix \(A^{T}A\):
\[
A^{T}A v_i=\sigma_i^2 v_i.
\]

So the singular values are
\[
\sigma_i=\sqrt{\lambda_i(A^{T}A)}.
\]

For each \(\sigma_i>0\), compute the matching left singular vector by
\[
u_i=\frac{Av_i}{\sigma_i}.
\]

If \(U\) or \(V\) is not complete, add orthonormal vectors to finish the basis.

\section{Geometry}

The SVD is the clean geometric order:
\[
x
\xrightarrow{\;V^{T}\;}
\text{rotate/reflection}
\xrightarrow{\;\Sigma\;}
\text{stretch}
\xrightarrow{\;U\;}
\text{rotate/reflection}.
\]

The singular values show the principal stretches of the matrix.

\section{Rank and the Four Subspaces}

If \(\sigma_1,\ldots,\sigma_r>0\) and the rest are zero, then \(r=\operatorname{rank}(A)\).

\begin{center}
\small
\begin{tabularx}{\linewidth}{@{}L{0.25\linewidth}Y@{}}
\toprule
Subspace & Basis from the SVD \\
\midrule
Row space \(C(A^{T})\) & \(v_1,\ldots,v_r\). \\
Nullspace \(N(A)\) & Remaining \(v_i\) with \(\sigma_i=0\). \\
Column space \(C(A)\) & \(u_1,\ldots,u_r\). \\
Left nullspace \(N(A^{T})\) & Remaining \(u_i\). \\
\bottomrule
\end{tabularx}
\end{center}

\section{SVD vs Eigenvalue Decomposition}

\begin{center}
\small
\begin{tabularx}{\linewidth}{@{}L{0.24\linewidth}Y Y@{}}
\toprule
Feature & Eigenvalue decomposition & SVD \\
\midrule
Basic form & \(A=S\Lambda S^{-1}\) & \(A=U\Sigma V^{T}\) \\
Works for & Diagonalizable square matrices & Every matrix \\
Numbers & Eigenvalues may be negative or complex & Singular values are nonnegative real numbers \\
Directions & Same space in and out & Input directions \(v_i\), output directions \(u_i\) \\
Best use & Dynamics and powers & Geometry, rank, least squares, approximation \\
\bottomrule
\end{tabularx}
\end{center}

\section{Symmetric Case}

For a symmetric positive definite matrix,
\[
A=Q\Lambda Q^{T}
\]
already has orthonormal eigenvectors and positive eigenvalues. In this special case, the SVD agrees with the eigenvalue decomposition.

For a symmetric matrix with negative eigenvalues, singular values use absolute sizes, not signs.

\section{Common Mistakes}

\begin{center}
\small
\begin{tabularx}{\linewidth}{@{}L{0.40\linewidth}Y@{}}
\toprule
Mistake & Correct exam habit \\
\midrule
Treating singular values as eigenvalues of \(A\). & Singular values come from \(A^{T}A\). \\
Allowing negative singular values. & Always take \(\sigma_i\ge 0\). \\
Using \(u_i=Av_i/\sigma_i\) when \(\sigma_i=0\). & Only divide by positive singular values. \\
Forgetting rectangular dimensions. & \(\Sigma\) has the same shape as \(A\). \\
Mixing up \(U\) and \(V\). & \(V\) belongs to the input side; \(U\) belongs to the output side. \\
\bottomrule
\end{tabularx}
\end{center}

\section{Final Exam Checklist}

\tightlabel{Checklist}
\begin{enumerate}
  \item Compute \(A^{T}A\).
  \item Find eigenvalues and orthonormal eigenvectors of \(A^{T}A\).
  \item Set \(\sigma_i=\sqrt{\lambda_i}\).
  \item Put the \(v_i\)'s into \(V\).
  \item Compute \(u_i=Av_i/\sigma_i\) for positive \(\sigma_i\).
  \item Complete orthonormal bases when needed.
  \item Write \(A=U\Sigma V^{T}\) with matching dimensions.
  \item Use zero singular values to identify nullspaces.
\end{enumerate}

\end{document}
